\documentclass[11pt,a4paper]{article}
\usepackage[margin=25mm]{geometry}
\usepackage{amsmath,amssymb,amsthm,hyperref}
\newtheorem{theorem}{Theorem}
\newtheorem{lemma}[theorem]{Lemma}
\newcommand{\gen}[1]{\langle#1\rangle}
\title{The pair partition in the paired equality case}
\author{Human Compression project: a component of the L044 recovery}
\date{}
\begin{document}\maketitle
Retain the setting of \emph{Neighbours of a square in a paired
selection}: a transitive string action on $n=2r>8$ vertices,
with two selected eligible edges of every colour and a connected
selection whose unique cycle is an alternating square. A square
is admissible if it is the unique cycle of such a selection.
We use the acquired four-root-tree lemma, the choice containing
all actual neighbours of a fixed square, the square transport
lemma, and the bound three for the actual degree of each vertex
of any admissible square. Eligibility always refers to the actual
omission group $G_i=\gen{\rho_j:j\ne i}$.

The source ladder strategy is due to Cameron--Fernandes--Leemans--Mixer,
\emph{Highest rank of a polytope for $A_n$}, arXiv:1605.09173v1,
Propositions~4.16--4.18. We give the selection and continuation
details explicitly. In particular, a slide below changes only
the rung edges: it does not silently replace an edge of another
colour farther along the selection.

\begin{lemma}[Squares with no common middle label]\label{automatic}
If an admissible square has colours $a,b$ with $|a-b|\ne2$,
every connected paired selection having that square as its
unique cycle already contains every actual edge incident with
the square.
\end{lemma}
\begin{proof}
Use the four root trees. An unselected actual $k$-edge from a
square vertex $p$ ends in $T_p$, and the tree path to its other
endpoint starts with a colour $\ell\ne k$. The path and the
unselected edge form a cycle. Eligibility forces both selected
$\ell$-edges onto that path. Transporting its first edge by any
square mark commuting with $\rho_\ell$ would require a selected
$\ell$-edge in a different root tree. Thus $\ell$ is adjacent
to both $a$ and $b$, forcing $|a-b|=2$, contrary to hypothesis.
\end{proof}

\begin{lemma}[Continuation along pairs]\label{continue}
Suppose the current selected square has rung colour $a$ and
longitudinal colour $c$, with $|a-c|\ge3$. Write its two
$a$-pairs as $B,F$, joined by the two $c$-edges. Suppose each
vertex of $B$ has a third actual edge of colour $d=c-\epsilon$,
where $\epsilon\in\{-1,1\}$ and these are the two selected
$d$-edges. Then every actual edge leaving the square at $F$
has colour $c+\epsilon$.

If the difference from $a$ is increasing in direction $\epsilon$,
such an edge has a parallel mate, both are selected, and moving
the square to these new endpoints changes only one selected
$a$-edge. All other selected edges are retained. Whenever all
previous vertices other than $F$ admit no selected exit, and
the new endpoints are outside those vertices, this extends the
selected ladder by one pair. If there is no exit at $F$, that
ladder is the whole connected selection.
\end{lemma}
\begin{proof}
Let $e$ be the colour of a further edge at a vertex of $F$.
If it commuted with $c$, transport by $\rho_c$ would give an
edge of colour $e$ at the corresponding vertex of $B$. That
vertex already has its three actual edges of colours $a,c,d$.
Since $e\notin\{a,c\}$ and $d$ is adjacent to $c$, none of
these can be the commuting colour $e$. Hence $e$ must be
$c-1$ or $c+1$. Lemma~\ref{automatic} makes this edge
selected. The two selected $d$-edges already meet $B$, so
$e\ne d$, leaving exactly $e=c+\epsilon$.

In the stated direction $e$ commutes with $a$. Its parallel
edge meets the other vertex of $F$ and is also selected by
Lemma~\ref{automatic}. The two new endpoints lie outside
the square, in its two different front root trees. Replace
the $a$-edge on $B$ by the transported $a$-edge on these
endpoints. The acquired square transport proof shows that
this leaves a connected paired selection with the new square
as its unique cycle. No other edge is changed. The final
assertions follow from connectedness: a selected path from
the displayed ladder to any other vertex would have to leave
through its current front.
\end{proof}

\begin{lemma}[A square whose labels differ by two]\label{gap}
There exists an admissible square with colours $a,a+2$.
\end{lemma}
\begin{proof}
If an admissible square has consecutive colours, choose a
selection containing all its actual neighbours. Connectedness
and $n>4$ give an edge leaving it. The root-tree lemma makes
its colour the outer neighbour of one square colour. It
commutes with the other, so square transport gives an
admissible square with label difference two.

Suppose, for contradiction, that no difference-two square
exists. Choose an admissible square with minimum label
difference, say colours $a<b$. The preceding paragraph
implies $b-a\ge3$. An edge leaving the square is adjacent
to a square colour. An inner attachment of colour $a+1$
or $b-1$ would give an admissible square of smaller difference.
Thus every attachment points outward. Reverse the full string
if necessary so that one has colour $c=a-1$.

Choose the selection containing all neighbours of this square.
The two parallel $c$-edges leave one $b$-pair, denoted $R_1$,
and reach a new pair $R_2$. Let $R_0$ be the other $b$-pair
of the original square. Both vertices of $R_1$ now have
actual degree three, with colours $b,a,c$. Any additional
edge at $R_0$ must fail to commute with $a$, since otherwise
transport to $R_1$ would exceed its degree. Its colour is
therefore $a-1$ or $a+1$. The former already has its two
selected edges at $R_1$, and all neighbours of the original
square are selected. The latter would give a smaller
admissible label difference. Thus $R_0$ has actual degree
exactly two.

Slide the square to $R_1,R_2$, with rung $b$ and longitudinal
colour $c=a-1$. Apply Lemma~\ref{continue} with
$\epsilon=-1$. Every successive extension is a parallel
pair of edges whose labels decrease by one. The rear pair
$R_0$ is closed; every subsequent interior vertex has its
three actual edges. New endpoints cannot be previously
constructed vertices, which already have their complete
degree or are in the closed rear pair. Thus the construction
continues along distinct pairs or stops at a closed front.
It must stop, since the labels cannot decrease indefinitely.
All full-graph exits have then been excluded, so transitivity
makes these pairs all of $\Omega$. But the only colours used
on them are $b,a,a-1,\ldots$; a colour strictly between $a$
and $b$ fixes every vertex. This is incompatible with having
two eligible edges of every colour. The contradiction proves
the lemma.
\end{proof}

\begin{lemma}[The two middle edges]\label{middle}
Choose a selection containing all actual neighbours of an
admissible square with colours $a,b=a+2$, and put $k=a+1$.
Both selected $k$-edges meet the square, at two adjacent
vertices. Their other endpoints have degree one in the
selected graph.
\end{lemma}
\begin{proof}
First suppose no $k$-edge meets the square. Some selected
edge leaves it, by connectedness. Its colour is $a-1$ or
$b+1$. Reverse labels if needed to choose $c=b+1$. It and
its parallel mate join one $a$-pair to a new pair. The
other $a$-pair of the square is closed: an additional edge
there must fail to commute with $b$, by transport to the
now saturated pair, so has colour $b-1=k$ or $b+1=c$.
The first is absent by hypothesis and the second already
uses its two selected edges at the saturated pair.

After sliding to the new square with colours $a,c$, its
label difference is three. Lemma~\ref{continue}, now with
$\epsilon=1$, continues along increasing longitudinal
labels. The closed rear and saturated interior vertices
force every new pair to be disjoint from the preceding ones.
The process ends because there are finitely many labels,
and the resulting full component omits colour $k$. As in
Lemma~\ref{gap}, transitivity gives a contradiction. Thus
some selected $k$-edge $pu$ leaves the square.

The vertex $u$ is a selected leaf. Indeed all selected
$a$- and $b$-edges already lie in the square, and the
involution $\rho_k$ already pairs $u$ with $p$. Any other
selected edge at $u$ would have a colour commuting with
$k$. Transport would give a fourth actual edge at $p$,
whose three edges have colours $a,b,k$. This is impossible.
The same reasoning applies to every selected middle edge
meeting the square.

Write $q=p^{\rho_a}$, $s=p^{\rho_b}$ and
$t=p^{\rho_a\rho_b}$. Suppose the second selected
$k$-edge does not meet the square. Since $u$ is a selected
leaf, connectedness forces an outer-colour attachment at
one of $q,s,t$. Reverse labels if necessary so that its
colour is $c=b+1$. As it commutes with $a$, and $p$ is
already saturated, its two parallel edges must leave $s,t$.
Both of these vertices are now saturated. The vertex $q$
is closed: an attachment there would be of colour $a-1$
or $b+1$, and commuting transport by $b$ or $a$,
respectively, would give a fourth edge at $t$ or $p$.
There is no actual $k$-edge at $q$ under the current
assumption, since all actual square neighbours are selected.

Slide the rung colour $a$ to the new pair. The new square
has label difference three. Continue outwards by
Lemma~\ref{continue}. Every slide changes only an
$a$-edge and retains $pu$ and the other selected $k$-edge.
The leaf $u$ remains a selected leaf: the selected
$a$-edges lie on the current square, the selected $b$-edges
are the original two longitudinal edges, and every other
possible colour except $k$ commutes with $k$ and fixes $p$,
hence fixes $u$. All previous nonleaf vertices are closed
or saturated. Each new longitudinal colour is at least
$b+1=k+2$, so it also fixes $u$ by this commutation.
Consequently new endpoints cannot be old vertices, and
any selected exit is through the current front.

At each front the only possible new colour is the next
increasing longitudinal colour. No second $k$-edge ever
enters the displayed ladder. When the process stops,
connectedness says this ladder and its leaf comprise
the entire selection, although it contains just one
selected $k$-edge. This is a contradiction.

The second middle edge therefore meets another square
vertex and ends in another selected leaf. It cannot meet
$t$, diagonally opposite $p$. In that case $p,t$ are
saturated and both leaves have selected degree one. Any
selected exit through $q$ or $s$ must have outer colour
$a-1$ or $b+1$; commuting transport to $p$ or $t$ again
contradicts degree three. The whole connected selection
would then have only six vertices. Since $n>8$, the two
middle edges instead meet adjacent square vertices.
\end{proof}

\begin{theorem}[The invariant pair partition]\label{pairs}
The group $G$ preserves a partition of $\Omega$ into $r$
pairs. In particular a primitive action of degree $n=2r>8$
cannot have the stated paired equality selection.
\end{theorem}
\begin{proof}
Use Lemmas~\ref{gap} and \ref{middle}. Reverse labels if
necessary so the middle edges of colour $k=a+1$ meet an
$a$-pair $R_0$ of the square of colours $a,b=a+2$.
Let $L$ be the pair of selected leaves, and let $R_1$ be
the other $a$-pair. The two vertices of $R_0$ have degree
three. Since $n>6$ and the leaves have selected degree one,
some selected edge must leave $R_1$. Its colour cannot
be $a-1$, which commutes with $b$ and would produce an
edge at the saturated pair $R_0$. It is therefore $b+1$.
Its parallel mate gives a new pair $R_2$.

Slide the $a$-rung to the new pair and continue by
Lemma~\ref{continue}. Both leaves remain selected leaves
by the same argument as in Lemma~\ref{middle}: each slide
changes only the rung colour, and the only possibly
noncommuting colours $a,b$ have their selected edges
elsewhere. All newly encountered colours commute with
$k$ and fix $R_0$, hence fix $L$. The remaining old
vertices are saturated. Thus the construction uses distinct
pairs and, by connectedness, ends only when it includes
all vertices of $Q$.

The selected colours in this completed graph are exactly
the consecutive interval $a,a+1,\ldots,j$, where $j$
is the last longitudinal colour. Every label
$0,\ldots,r-1$ has two selected edges, so $a=0$ and
$j=r-1$. We obtain pairs
\[
                        L,R_0,R_1,\ldots,R_{r-2}.
\]
There are two edges of colour $1$ between $L$ and $R_0$,
and two edges of colour $i$ between $R_{i-2}$ and
$R_{i-1}$ for $2\le i\le r-1$. Each $R_j$ is an
actual $0$-edge, as supplied by successive commuting
transports. Both vertices of every $R_j$ with $j<r-2$
have their three actual edges.

The terminal pair $R_{r-2}$ has exactly its actual edges
of colours $0,r-1$. Indeed apply the first paragraph of
Lemma~\ref{continue} to the final square: any additional
edge would have colour $r-2$ or $r$. The latter is outside
the label set. The former already has its two selected
edges incident with the preceding pair, and all actual
neighbours of the final square are selected, by
Lemma~\ref{automatic}. Here $r\ge5$, so its label
difference is at least four, as required.

There are no vertices left outside these pairs. Since
every vertex other than the two in $L$ has all its actual
edges accounted for, the only possible additional actual
edge joins the two leaves. Such an edge must have colour
$0$: a colour different from $0$ would close the path
from one leaf to the other through its two $1$-edges and
the original $0$-rung on $R_0$, making a cycle with just
one eligible $0$-edge. The original rung is eligible
whether or not it is retained in the final selection.
Thus $\rho_0$ interchanges each pair $R_j$ and either
fixes or interchanges the two leaves. Each $\rho_i$,
$i\ge1$, swaps the corresponding two adjacent pairs by
its displayed two edges and fixes the remaining vertices.
Every generator therefore preserves the stated partition.
\end{proof}

\paragraph{Exact scope.}
The acquired normal-form and equality-connectedness results,
together with Theorem~\ref{pairs}, exclude $2r=n$ for
primitive actions with a paired selection and $n>8$.
The paired-cut count already gives $2r\le n$, so in a
primitive action of degree $n\ge12$ it gives
$r\le\lfloor(n-1)/2\rfloor$. This recovers the needed
all-paired branch of L044. It does not classify every
imprimitive equality group or supply the separate general
transitive-imprimitive omission branch, induction bases,
small maxima or lower existence. Those remain part of the
complete principal-source contract. No whole-source or new
open-conjecture credit is asserted here.
\end{document}
